\documentclass[11pt]{article}
\usepackage{amsmath,amssymb,hyperref}
\title{Forecasts, preparation, and clearance around unfinished work}
\date{24 November 2026}
\begin{document}
\maketitle

In a finite maintenance model, the best use of a warning changes at an exact cost boundary. An early warning supports preparing in a spare worker slot, or diverting when the slot is already owed to old work. At preparation cost $3/8$ and diversion cost $3/4$, a controller that can choose either attains worst expected service cost $13/32$. The best preparation-only policy costs $14/32$ and the best diversion-only policy costs $16/32$. These costs apply to the declared route and uncertainty model. The clearance decision comes from an existing scheduler for real obligations in the model, not from the forecast bit.

\section{What arrives before each decision}
The old schedule is the two-handoff replay of \href{https://github.com/eternal-november/eos-research-programmes/pull/353}{\#353}. Before the mission begins, the old desk has admitted two totes, the first due for closure at tick 2 and the second at tick 3. The relay remains required at every tick through 3. At tick 2 the maintenance plant may have a delayed B health report or a failed B; those worlds are not distinguished when remaining work is assigned. The inherited closures and relay are mandatory in both. The certified intake inspection of the first tote is a private input to the old-work dispatcher and authorization gate. The mission controller cannot read that observation, the old live roster, probe outcomes, or the physical tape. It sees only appointed mission receipts at the indicated cuts. Lower lead times are not equivalent to a change in physical clearance.

We replace the optional new-tote slot query of \#353 by a shielding installation, using exactly its current-gate authority and worker-time requirement $C_2$. No extra tote is admitted. For an intact intake, assign the old closure to $A_2$ and the relay to $D_2$, leaving $C_2$ for this installation. For a damaged intake, the first closure, tick-2 relay and installation have only $D_2,C_2$ in the intersection graph, so the joint slot assignment fails. This is a genuine Hall deficit when B failed; a later proof that B was merely delayed would allow a relay reassignment in that completion. It would not license that schedule on the unresolved history. The relay and the second closure have a separate proof at tick 3: they use $C_3$ and $D_3$ respectively. The replay calls the exact original matching and timeline routines, then independently checks feasible assignments by enumeration.

The mission controller can choose to prepare at tick 2, when an installation is permitted and it sees an early clearance; installation changes the arrangement from tick 3. If it has not prepared, it can choose a detour at tick 3. This requires a separate E technician and protected switch slot $E_3$; the synthetic plant gives E no old duties and excludes interchange with C or D. A separate control instead requires $C_3$ for switching. That slot must carry the relay because the second closure takes $D_3$, so no detour is then feasible. This changes the shared capacity, not the forecast law or installation certificate.

Every route enters its zone at tick 4 and completes at tick 6 or 7. A later disturbance on the normal zone at tick 5 causes a one-unit safe reactive hold without preparation. The same zone with shielding absorbs the interruption and causes no hold. The alternative route crosses a different, unaffected zone after the independent switch; its control never receives a normal-zone outcome as its own. Preparation and diversion have service costs $c,d$ respectively, in the unit of that reactive hold. The disturbance potential, exposed zone and shielding effect are stipulated transition rules, not estimated causal effects. All evaluated branches have exclusive slots, old closure and relay intact, and mission completion by 7. The optimized loss scores impact within that deadline, not risk of injury or of missing the final deadline.

An outcome is determined by a common exogenous tape chosen before actions: intact or damaged first intake, warning $W$, potential later normal-zone disturbance $H$, and B's hidden condition. A single tape allows separate route executions with the same exogenous values. Actual route outcomes come from crossing the chosen zone under its arrangement; no estimator retrospectively treats a baseline observation as an observation on a route not taken. Forecast and disturbance may covary with old workload. B's condition is an independent fair bit under both compared laws and is not used as decision evidence.

\section{Matched forecast laws and the parameter boundary}
Let $F$ denote existence of an authorized spare $C_2$ slot on the unresolved history, derived above. In order $(H,W)=(1,1),(0,1),(1,0),(0,0)$, both laws have counts $(3,1,1,3)$ out of eight. Among the counts, the intact intake counts are $(2,0,0,2)$ in A and $(1,1,1,1)$ in B. The plant rule yields the corresponding $F=1$ counts. This assigns weights to the same physically possible tapes, with some zero weights. The $(H,W)$, $(H,F)$ and $(W,F)$ distributions agree, and all four $(F,W)$ cells have mass $1/4$. The families differ at higher order; multiplying the warning model by separate clearance frequencies cannot recover the optimal decision. In fact the two displayed laws are endpoints within distributions on $(F,W,H)$ with those pairwise marginals: the $F=1$ counts are $(t,2-t,2-t,t)$ with $1\le t\le 2$. Each randomized policy's loss is affine in $t$. For this calculation, the plant, cost model and B independence remain fixed.

Consider $0\le c\le 1/2$ and $1/2\le d\le1$. In cell $(F,W)=(1,1)$ preparation weakly dominates both alternatives under both laws; in $(0,0)$ staying weakly dominates. In clear silence, diversion is dominated by preparation, and in blocked warning preparation is inadmissible. Let $q$ be the probability of preparing in clear silence and $s$ the probability of diverting in blocked warning. In units of four times expected cost the two endpoint losses are
\[
 L_A=c+1+cq+(d-\tfrac12)s,\qquad
 L_B=c+\tfrac32+(c-\tfrac12)q+(d-1)s.
\]
The first is nondecreasing, the second nonincreasing in each parameter, with $L_A-L_B=(q+s-1)/2$. Moving within the square to $q+s=1$ does not increase the maximum; there, assign the unit weight to the cheaper coefficient in $L_A$. Hence the exact randomized minimax value is
\[
 V=\frac{1+c+\min(c,d-\tfrac12)}4.
\]
This includes all randomized causal policies here: the only meaningful choice before disruption is the arrangement determined at the declared cuts; in the early regime both bits are available then, and an independent random seed gives a mixture of deterministic action tables. The rule prepares whenever clear if $c<d-1/2$. If $c>d-1/2$, it prepares on clear warning, diverts on blocked warning, and stays on silence. Either rule attains the value at equality, with possible additional ties on the parameter edges.

Without diversion the minimax value is $(1+2c)/4$, attained by preparing whenever clear. Without preparation the equal mixture of A and B gives disturbance posterior $3/4$ in warning and $1/4$ in silence at both clearance grades. An optimal warning-only switch or stay attains that mixture bound under both laws. Thus the exact value without preparation is $(1/2+2\min(3/4,d))/4$. These are optimized restricted policies, not straw comparisons against always diverting or ignoring all evidence. A cost savings from diversion requires the strict region $c>d-1/2$ as well as actual access to a diversion resource. The $C_3$ collision removes precisely that action and restores the preparation-only value without changing the warning or the installation clearance.

\section{Separating availability, delivery, and authority}
For the late-clearance control, keep every physical tape, the gate's private intake knowledge, and the reported clearance contents fixed. Only the mission receipt moves from tick 2 to tick 3. A preparatory submission without a certificate is not allowed by the incumbent rule, and all warning values still have feasible and infeasible physical completions. Some tapes permit installation physically, yet no uniformly safe unconditional early installation can be authorized from the controller's view. The late value is the exact no-preparation value. The installed certificate consumption uses the same distinction between evidence contents, arrival and authorized interval as \href{https://github.com/eternal-november/eos-research-programmes/pull/354}{\#354}; it does not invoke that study's lane alphabet bound.

A different interface can restore the early value on those same tapes. An atomic conditional authorization is sent at tick 2 using only the warning and a configured rule: ask the gate to install if its independent matching witness allows it. The gate grants only with a sound slot witness and otherwise leaves the old plan intact. At tick 3 the mission controller has the delayed receipt, can distinguish grant from no installation, and can still select a feasible detour. A rejected request consumes no worker-time and its observable outcome arrives before the later cut in this control. With these extra authority and feedback assumptions, every early action table can be implemented: submit when the clear branch would prepare for that warning, and execute its late non-prepare branch or the blocked branch if it did not grant. The request is not a fictitious early observation of $F$. If refusal itself consumed a critical slot or release occurred after the switch deadline, preservation would need a new proof.

At $(c,d)=(3/8,3/4)$, keeping clearance early but delivering $W$ only at tick 3 yields $14/32$, even when $W$ can still govern diversion. To see the lower bound, let $p$ be preparation probability when clear, which cannot depend on $W$ at tick 2. Under the equal endpoint mixture, conditional disturbance probabilities by warning are $3/4$ and $1/4$. The optimal unprepared continuation is worth $1/2$ averaged over fair warnings; the blocked regime also costs $1/2$. Hence the mixture bound is $1/2-p/16$, minimized at $p=1$. Preparing whenever clear and always staying otherwise attains $14/32$ under both laws. Delivery at tick 4 cannot improve that feasible policy and has the same value. Thus late forecast and late clearance remove different options even on the same tapes.

\section{What was reused}
The replay preserves \#353's old two admissions, old closure checks, authority handoff, slot relation and continuous relay proof. One new application of the existing admission query derives the installation condition. The alternative slot proof, timing paths, route exposure equations, matched tape laws and cost optimization are new; they do not follow from the slot proof. \#354 guides the observation and authorization specification, but its proof about held lane intervals is not imported as a guarantee for these worker-time slots. A factored direct proof uses precisely the same slot witnesses, timed receipt rule and route checks. No reduction relative to that baseline is shown. The comparison is with existing work on \href{https://arxiv.org/abs/2103.05526}{distributed predictive control with information-age forecasts} and \href{https://doi.org/10.1137/20M1370227}{two-stage distributionally robust optimization}; this finite timing and capacity calculation claims neither a new general optimization method nor prediction novelty.

An exact independent deterministic-table and two-policy-mixture search matches the closed formula at 20 cost pairs. The executable trace checker covers 16 matched exogenous tapes, 12 source scheduling queries with independent assignment enumeration, and 304 feasible action interventions across timing and resource modes. The source replay is SHA-256 pinned and checked at import. These are synthetic checks and a scoped argument, not physical testing, Lean certification, full C7 acceptance, or a frozen C11 transfer.
\end{document}
